% ============================================================
% Appendix
%
% Supplementary recursive synchronization material for readers
% who want deeper operational or architectural detail.
%
% Structure:
%   A — Example Recursive Workflow Lifecycle
%   B — Example Specification Structures
%   C — Example Repository Architecture
%   D — Publication Workflow Example
%   E — Reflective Audit Examples
% ============================================================

\section{Example Recursive Workflow Lifecycle}
\label{sec:appendix-workflow}

The following walkthrough traces a single end-to-end recursive workflow cycle as it would be executed under the \reflector{} architecture described in \Cref{sec:reflector-framework}. The purpose is to make the lifecycle concrete: each stage is named, its inputs and outputs are described, and its relationship to synchronization and governance is made explicit. This example is deliberately lightweight so that the structural pattern remains visible without being obscured by domain-specific detail.

Canonical companion artifacts for these appendix examples are provided under \texttt{paper/examples/canonical/}, including synchronization and publication specifications, recursive workflow manifests, reflective audit examples, and lightweight synchronization diagrams.

\subsection{Canonical Lifecycle Stages}

A complete \reflector{} workflow cycle proceeds through eight ordered stages:

\begin{center}
Human Goal
$\rightarrow$ Scoped Issue
$\rightarrow$ Specification Alignment
$\rightarrow$ AI-Assisted Execution
$\rightarrow$ Generated Artifact
$\rightarrow$ Reflective Audit
$\rightarrow$ Synchronization Checkpoint
$\rightarrow$ Recursive Continuation
\end{center}

Each stage is a bounded transformation with defined inputs, expected outputs, and explicit acceptance conditions.

\subsection{Stage Descriptions}

\begin{itemize}
  \item \textbf{Human Goal.} A human operator identifies an objective and expresses it in natural language. The goal is intentionally high-level and does not prescribe implementation. At this stage, the key governance decision is whether the goal is ready to be scoped---whether its intent is clear enough to produce a bounded, auditable issue.

  \item \textbf{Scoped Issue.} The goal is decomposed into a repository issue with explicit acceptance criteria, bounded scope, and a link to the relevant specification. Issue scoping transforms broad intent into a delegable unit of work. Acceptance criteria serve as the reference surface for the reflective audit that follows.

  \item \textbf{Specification Alignment.} Before delegation begins, the relevant specification artifact is consulted or updated to confirm that the issue is consistent with declared architectural intent. If no specification exists for the relevant domain, one is drafted at this stage. This step prevents execution from drifting away from architectural commitments at the outset \cite{Parnas1972Modularization}.

  \item \textbf{AI-Assisted Execution.} An AI system receives the scoped issue and specification context, then executes within those boundaries to produce candidate artifacts. The scope boundary is explicit: the AI is authorized to modify the artifact surface declared in the issue and no more. Execution may be iterative within the delegated scope.

  \item \textbf{Generated Artifact.} Execution produces one or more repository-committed artifacts: code changes, documentation sections, figures, test files, or specification updates. The artifact is committed with enough attribution context---commit message, issue reference, and change summary---that its provenance can be reconstructed without relying on ephemeral session state \cite{Lamport1978TimeClocks}.

  \item \textbf{Reflective Audit.} The generated artifact is compared against the issue's acceptance criteria and the specification anchor. The audit asks: Does the artifact address the declared scope? Does it preserve cross-artifact consistency? Does it introduce drift signals? If the audit produces concerns, the workflow branches: the artifact is returned for correction before the checkpoint is attempted.

  \item \textbf{Synchronization Checkpoint.} If the reflective audit is satisfied, a synchronization checkpoint is recorded. The checkpoint documents the artifact state, the audit evidence, and the continuation authorization. This creates an explicit, inspectable record of the moment when the recursive workflow was cleared to proceed \cite{Chandy1985DistributedSnapshots}.

  \item \textbf{Recursive Continuation.} With the checkpoint closed, the next scoped issue is authorized. The recursive workflow continues from a refreshed and reconciled state rather than from accumulated ambiguity. Over many cycles, these checkpoints compose into a coherent artifact history that supports both incremental progress and retrospective audit.
\end{itemize}

\subsection{Recursive Issue Orchestration}

Recursive orchestration means that the output of one lifecycle cycle becomes the input context for the next. The ``Human Goal'' of a later cycle is typically scoped by the artifact state left at the previous checkpoint. This composition is what makes the lifecycle recursive rather than merely iterative: each cycle inherits the synchronized state of its predecessors, and drift from one cycle is constrained from propagating into the next by the intervening synchronization checkpoint.

Well-managed recursive orchestration exhibits several properties \cite{Wiener1948Cybernetics,Beer1972BrainFirm}:
\begin{itemize}
  \item \textbf{Bounded depth:} no recursive cycle proceeds indefinitely without a synchronization event.
  \item \textbf{Inspectable history:} every cycle leaves a committed artifact trail that can be re-entered at any prior checkpoint.
  \item \textbf{Stabilizing refinement:} repeated cycles converge toward coherent artifact state rather than accumulating hidden inconsistency.
  \item \textbf{Governed continuation:} human authority remains coupled to the moments that authorize recursive progression.
\end{itemize}

% ============================================================

\section{Example Specification Structures}
\label{sec:appendix-specs}

Specifications function as synchronization anchors throughout the \reflector{} workflow. The following examples illustrate what representative specification artifacts look like in a \reflector{}-aligned repository. These are structural sketches rather than prescriptive schemas; the exact format should match the repository's existing conventions and documentation tooling \cite{Parnas1972Modularization}.

\subsection{Specification Taxonomy}

A \reflector{}-aligned repository typically organizes specifications by functional domain:

\begin{lstlisting}[language=bash, caption={Example specification directory structure.}, label={lst:spec-dir}]
specs/
  synchronization.spec.md
  publication.spec.md
  visual-synapse.spec.md
  workflow.spec.md
  publication/
    publication-manifest.spec.md
    publication-workflow.spec.md
    semantic-content.spec.md
    renderer-architecture.spec.md
\end{lstlisting}

Each specification file declares the scope, constraints, evidence requirements, and continuation conditions for a specific domain. Together they form the specification layer described in \Cref{sec:implementation-examples}.

\subsection{Synchronization Specification}

A synchronization specification records the checkpoint structure, evidence requirements, and continuation criteria for a recursive workflow domain. A representative structure includes:

\begin{lstlisting}[caption={Synchronization specification structure (Markdown front matter and body).}, label={lst:sync-spec}]
# Synchronization Specification

## Scope
Defines checkpoint conditions for recursive issue orchestration.

## Checkpoint Types
- specification-checkpoint: re-align intent before execution
- artifact-checkpoint:      validate generated artifact set
- architectural-checkpoint: reconcile against system boundaries
- milestone-checkpoint:     evaluate readiness for continuation

## Evidence Requirements
Each checkpoint must record:
- artifact state at checkpoint time
- audit result (pass / conditional / fail)
- continuation authorization (human sign-off)

## Continuation Criteria
Recursive continuation is authorized only when:
- all declared acceptance criteria are satisfied
- no unresolved drift signals remain open
- artifact set is self-consistent
\end{lstlisting}

\subsection{Publication Specification}

A publication specification declares the semantic content structure, renderer targets, metadata requirements, and manifest format for a publication pipeline. A representative structure includes:

\begin{lstlisting}[caption={Publication specification structure.}, label={lst:pub-spec}]
# Publication Specification

## Semantic Content
- Sections: structured LaTeX source files in sections/
- Figures:  PNG exports in figures/ with manifest.md
- References: BibTeX entries in references.bib

## Metadata
- Title, author, version: macros/metadata.tex
- Publication status: \paperstatus macro

## Renderer Targets
- arxiv:   single-file LaTeX + bibliography bundle
- ieee:    IEEEtran class with double-column layout
- acm:     acmart class with rights statement block
- pages:   HTML export via pandoc or LaTeX4ht

## Manifest Requirements
Each release must include:
- compiled PDF with stable figure references
- source archive reproducible from repository state
- artifact version and build timestamp
\end{lstlisting}

\subsection{Workflow Specification}

A workflow specification enumerates the phases of a recursive execution cycle and the conditions governing each transition:

\begin{lstlisting}[caption={Workflow specification structure.}, label={lst:workflow-spec}]
# Workflow Specification

## Phases
1. goal-identification:    human articulates intent
2. issue-scoping:          decompose into bounded delegable unit
3. specification-alignment: confirm architectural consistency
4. execution:              AI-assisted artifact generation
5. reflective-audit:       check artifact against acceptance criteria
6. synchronization:        record checkpoint and authorize continuation
7. recursive-continuation: open next scoped issue

## Phase Transition Conditions
- goal -> issue:    goal is unambiguous and ready to scope
- issue -> spec:    acceptance criteria are enumerated
- spec -> execute:  relevant specifications are current
- execute -> audit: artifacts are committed to repository
- audit -> sync:    audit result is pass or conditional-pass
- sync -> continue: human governance authorization recorded
\end{lstlisting}

\subsection{Specifications as Synchronization Anchors}

The common thread across all specification types is that they establish a stable reference before execution and persist as the authoritative ground truth after it. This is what makes them synchronization anchors: they do not move with the workflow; instead the workflow is checked against them. Recursive execution without stable specifications is analogous to navigation without a fixed reference point---local movement may be rapid and confident while global position becomes increasingly uncertain \cite{Hutchins1995CognitionWild}.

% ============================================================

\section{Example Repository Architecture}
\label{sec:appendix-repository}

This appendix section illustrates what a \reflector{}-aligned repository looks like structurally. The layout reflects the semantic/render separation, recursive artifact organization, and publication manifest conventions described in \Cref{sec:implementation-examples}.

\subsection{Reference Repository Layout}

\begin{lstlisting}[language=bash, caption={Reference repository architecture for a \reflector{}-aligned project.}, label={lst:repo-arch}]
paper/
  paper.tex            # Thin orchestration wrapper
  references.bib       # BibTeX bibliography
  macros/
    metadata.tex       # Paper metadata commands
  styles/
    reflector.sty      # Publication style (swap to change renderer)
  sections/            # Semantic content (LaTeX section files)
  figures/             # PNG exports and figures manifest
    manifest.md
  diagrams/            # Source diagram files (Excalidraw, etc.)
  assets/              # Static assets (logos, images)
  publication/         # Publication manifests and release bundles

specs/
  synchronization.spec.md
  publication.spec.md
  workflow.spec.md
  publication/
    publication-manifest.spec.md
    publication-workflow.spec.md
    renderer-architecture.spec.md

workflows/             # Recursive orchestration workflow definitions
scripts/               # Build, validation, and diagnostic scripts
.github/
  workflows/           # CI workflows (build, pages, release)
\end{lstlisting}

\subsection{Architectural Principles}

Several principles drive this layout:

\begin{itemize}
  \item \textbf{Semantic/render separation.} The \texttt{sections/} directory holds semantic content only---argumentation, structure, and figure references. All typographic and layout decisions are delegated to the style layer (\texttt{styles/reflector.sty}). Swapping the style file targets a different publication format without touching section content.

  \item \textbf{Metadata centralization.} Document identity, authorship, version, and status are declared once in \texttt{macros/metadata.tex} and consumed everywhere via macros. This prevents the common drift failure mode in which metadata fields diverge across a document as it evolves.

  \item \textbf{Recursive artifact organization.} Specifications, workflows, diagrams, and publication configuration each occupy distinct directories. This makes artifact relationships navigable and auditable without requiring knowledge of which tool produced which file.

  \item \textbf{Publication manifests.} The \texttt{publication/} directory holds manifest files that declare exactly which inputs should be assembled for a given release target. A manifest-driven build is deterministic: the same manifest and repository state produce the same output, regardless of when the build is run \cite{Vogels2009EventuallyConsistent}.

  \item \textbf{CI as recursive validation.} GitHub Actions workflows in \texttt{.github/workflows/} function as automated participants in the synchronization checkpoint loop. Each pull request triggers a build that validates the artifact set before integration; each merge represents an authorized synchronization event.
\end{itemize}

\subsection{Synchronization-Aware Repository Conventions}

A synchronization-aware repository observes naming and organization conventions that lower the cognitive cost of reflective audit \cite{Simon1957ModelsMan}:
\begin{itemize}
  \item Issue titles encode scope explicitly (\textit{e.g.,} \texttt{[paper] Write synchronization section}).
  \item Commit messages reference the authorizing issue and summarize the artifact surface changed.
  \item Pull request descriptions enumerate the acceptance criteria being satisfied and the specification anchors consulted.
  \item Labels and milestones annotate which recursive cycle each issue belongs to, making workflow progress legible at a glance.
\end{itemize}

These conventions do not require new tooling; they require only that teams treat repository metadata as a first-class synchronization artifact rather than as administrative overhead.

% ============================================================

\section{Publication Workflow Example}
\label{sec:appendix-publication}

This appendix section traces a complete publication workflow as it would operate under the \reflector{} architecture, from initial semantic content through final renderer-specific artifact. The workflow reinforces the publication architecture principles discussed in \Cref{sec:implementation-examples} with a concrete, stage-by-stage example.

\subsection{Publication Pipeline Stages}

The canonical \reflector{} publication pipeline proceeds as follows:

\begin{center}
Semantic Content
$\rightarrow$ Publication Manifest
$\rightarrow$ Renderer Selection
$\rightarrow$ Build and Validation
$\rightarrow$ Generated Artifact
$\rightarrow$ Publication Checkpoint
$\rightarrow$ Distribution
\end{center}

Each stage is a bounded transformation that can be audited, versioned, and re-executed independently from repository state.

\subsection{Stage Descriptions}

\begin{itemize}
  \item \textbf{Semantic Content.} Authors produce structured source files in the \texttt{sections/} directory. Each file contains argumentation, figure references, and citation keys, but no typographic or layout decisions. The bibliography, figures, and metadata are maintained in their respective directories. At this stage, content correctness is evaluated against the publication specification \cite{Parnas1972Modularization}.

  \item \textbf{Publication Manifest.} A manifest file declares which content, metadata, and style artifacts are assembled for a given release. The manifest specifies the renderer target, the section inclusion order, the figure resolution requirements, the bibliography backend, and the expected output filename. Manifest-driven assembly ensures that publication builds are reproducible from repository state alone.

  \item \textbf{Renderer Selection.} The style layer is selected based on the renderer target declared in the manifest. Renderer targets in common use include:
  \begin{itemize}
    \item \texttt{arxiv} --- single-file \LaTeX\ bundle with \texttt{article} class, suitable for arXiv submission.
    \item \texttt{ieee} --- \texttt{IEEEtran} class with double-column layout, for IEEE conference and journal submission.
    \item \texttt{acm} --- \texttt{acmart} class with rights statement and CCS concepts, for ACM venues.
    \item \texttt{pages} --- HTML export pipeline for GitHub Pages or documentation hosting.
    \item \texttt{internal} --- unrestricted layout for preprint, review, or working-draft distribution.
  \end{itemize}
  Renderer selection is a style-layer swap; it does not touch the semantic content \cite{Vogels2009EventuallyConsistent}.

  \item \textbf{Build and Validation.} A build script assembles the inputs declared in the manifest, compiles the document using the selected renderer, and validates the output: bibliography completeness, figure reference resolution, cross-reference consistency, and metadata correctness. Validation failures halt the pipeline; all errors are reported with diagnostic detail before the artifact is accepted.

  \item \textbf{Generated Artifact.} A validated, renderer-specific artifact---typically a PDF and a source archive---is committed to the \texttt{publication/} directory and tagged with version and build metadata. The artifact is accompanied by a build log that documents the exact inputs and configuration used to produce it, supporting deterministic reproduction.

  \item \textbf{Publication Checkpoint.} A synchronization checkpoint records the artifact state, the build validation evidence, and the governance authorization for release. This checkpoint mirrors the structure of the recursive workflow checkpoint described in \Cref{sec:appendix-workflow}, but scoped to the publication domain. Only artifacts that pass the publication checkpoint are eligible for distribution.

  \item \textbf{Distribution.} The validated artifact is distributed to the declared targets: uploaded to arXiv, submitted to a venue, published to GitHub Pages, or circulated as a preprint. Distribution is a one-way transition: the distributed artifact is archived and the relevant publication specification is updated to record the release.
\end{itemize}

\subsection{Multi-Target Publication}

A \reflector{}-aligned publication system can target multiple renderers from the same semantic content without modification. The following example illustrates a repository that simultaneously maintains artifacts for three targets:

\begin{lstlisting}[language=bash, caption={Multi-target publication artifact layout.}, label={lst:multi-target}]
publication/
  arxiv/
    paper.tex          # flattened single-file bundle
    paper.bbl          # precompiled bibliography
    figures/           # figure exports at submission resolution
  ieee/
    paper.pdf          # compiled with IEEEtran style
  pages/
    index.html         # HTML export for web hosting
  manifests/
    arxiv.manifest.md
    ieee.manifest.md
    pages.manifest.md
\end{lstlisting}

Each manifest declares its own renderer target, figure requirements, and metadata fields. The semantic content in \texttt{sections/} is shared across all three. This separation means that a correction to section text propagates to all renderer targets on the next build cycle without requiring manual synchronization \cite{Chandy1985DistributedSnapshots}.

\subsection{Publication as Synchronization Infrastructure}

The publication workflow above is not merely an output delivery mechanism; it is synchronization infrastructure. The manifest declares what the publication artifact is supposed to contain; the build validates whether the current content state produces that artifact correctly; the checkpoint authorizes distribution only when the full pipeline passes. This mirrors the recursive workflow structure throughout the paper: intent is declared, execution is bounded, artifacts are audited, and continuation is governed. Publication orchestration therefore exemplifies \reflector{}'s broader claim that synchronization is a systems requirement, not an administrative afterthought.

% ============================================================

\section{Reflective Audit Examples}
\label{sec:appendix-audit}

Reflective audit is the observability function of the \reflector{} architecture, described in detail in \Cref{sec:reflective-auditing}. The following examples illustrate how reflective audit operates in four concrete engineering contexts: pull request review, specification reconciliation, figure synchronization, and publication validation. Each example identifies what is being audited, what evidence is consulted, and what an audit finding looks like in practice.

\subsection{Pull Request Audit}

A pull request is a synchronization boundary in which a proposed artifact change is exposed for review before integration \cite{Jimenez2024SWEBench}. A \reflector{}-aligned pull request audit asks four questions:

\begin{enumerate}
  \item \textbf{Scope conformance.} Does the diff address only the artifact surface declared in the authorizing issue? Changes outside the declared scope are flagged as boundary violations, even if individually correct.
  \item \textbf{Acceptance criteria coverage.} Does the artifact change satisfy the acceptance criteria enumerated in the issue? Each criterion is checked explicitly; partial satisfaction is noted for the next recursive cycle.
  \item \textbf{Specification consistency.} Does the change remain consistent with the specification anchor declared in the issue? If implementation diverges from the specification, either the implementation or the specification must be updated before the checkpoint is closed.
  \item \textbf{Cross-artifact integrity.} Do related artifacts---tests, documentation, figures, configuration files---remain self-consistent after the proposed change? A code change that invalidates existing test coverage or leaves documentation stale is a synchronization failure even if the code itself is locally correct.
\end{enumerate}

A pull request audit produces one of three outcomes: \textit{approved} (all criteria satisfied, checkpoint authorized), \textit{conditional} (minor issues noted, continuation permitted with deferred cleanup issue), or \textit{blocked} (significant drift detected, artifact returned for correction before checkpoint).

\subsection{Specification Reconciliation}

Specifications evolve over time as architectural understanding deepens. Specification reconciliation is the audit process that ensures that the specification layer remains consistent with the current implementation and with other specifications in the repository.

A specification reconciliation audit compares:
\begin{itemize}
  \item The declared scope of the specification against the current implementation surface.
  \item The evidence requirements in the specification against the evidence currently produced by the workflow.
  \item The cross-references between specifications against each other to detect inconsistencies introduced by independent updates.
\end{itemize}

When a reconciliation audit finds divergence---for example, a synchronization specification that declares a checkpoint structure no longer supported by the workflow tooling---the finding is recorded as a scoped issue with an explicit correction scope and a link to the affected specification. Recursive continuation in the affected domain is deferred until the reconciliation is complete \cite{Beer1972BrainFirm}.

\subsection{Figure Synchronization}

Figures are a particularly common source of cross-artifact drift in technical publications. Figure synchronization audits verify that the figure set remains consistent with the document in four respects:

\begin{itemize}
  \item \textbf{Reference existence.} Every \texttt{\textbackslash{}label} referenced via \texttt{\textbackslash{}Cref} in the text corresponds to an actual figure environment in the document; no dangling references remain.
  \item \textbf{Caption alignment.} Figure captions accurately describe the content they accompany. A placeholder caption carried over from an earlier draft and left unchanged after figure content is updated is a synchronization failure.
  \item \textbf{File availability.} Every figure file referenced in the document exists at the declared path, at an acceptable resolution, and in a format supported by the build pipeline.
  \item \textbf{Manifest consistency.} The figures manifest (\texttt{figures/manifest.md}) records all canonical figure assets; the manifest is checked against the document's figure references to confirm that no figure is referenced without a manifest entry and no manifest entry is orphaned.
\end{itemize}

Figure synchronization audits are particularly valuable at publication checkpoints, where a single missing or mis-captioned figure can invalidate the entire submitted artifact.

\subsection{Publication Validation}

A publication validation audit is a comprehensive reflective review run before any distribution event. It combines the concerns of pull request audit, specification reconciliation, and figure synchronization with additional publication-specific checks:

\begin{itemize}
  \item \textbf{Bibliography completeness.} Every citation key referenced in the text resolves to a bibliography entry; no entries are missing or malformed.
  \item \textbf{Metadata correctness.} Title, author, version, date, and status fields are current and consistent across the metadata layer (\texttt{macros/metadata.tex}), the compiled PDF properties, and any publication manifest files.
  \item \textbf{Build reproducibility.} The manuscript compiles without warnings or errors from the current repository state. Build reproducibility is confirmed by running a clean build from a fresh working directory rather than relying on cached intermediate files.
  \item \textbf{Renderer-target conformance.} The compiled artifact satisfies the submission requirements of the declared renderer target: page limits, column format, margin constraints, and any venue-specific packaging requirements.
  \item \textbf{Artifact provenance.} The final artifact is traceable to specific repository commits and specification versions, so that the publication checkpoint can be re-examined without reconstructing the conversational context of the development session \cite{Lamport1978TimeClocks}.
\end{itemize}

When all checks pass, the publication validation audit produces a signed checkpoint record that authorizes distribution. When any check fails, the specific failure is documented as a scoped correction issue and the distribution event is deferred. This structure ensures that publication quality is governed by explicit synchronization evidence rather than by informal confidence.

\subsection{Recursive Repository Audit}

Over extended recursive workflows, the repository itself can accumulate drift: orphaned specification files, stale workflow definitions, obsolete issue references, and figure files that are no longer referenced by any document section. A periodic recursive repository audit surfaces this accumulated drift before it compounds into structural incoherence.

A recursive repository audit examines the repository holistically \cite{Wiener1948Cybernetics}:
\begin{itemize}
  \item All open issues are cross-referenced against their declared milestone and specification anchor to confirm that none have drifted outside their authorized scope.
  \item All specification files are checked for internal consistency and cross-specification coherence.
  \item All figures and assets are validated against the figures manifest and against current document references.
  \item The CI workflow configuration is confirmed to reflect the current build and validation requirements.
  \item The publication manifests are confirmed to reflect the current renderer targets and artifact organization.
\end{itemize}

A recursive repository audit is a high-cost operation and is therefore not run at every cycle. It is most valuable at major milestone boundaries---before a public release, at the end of a development phase, or when significant architectural changes have accumulated across many smaller cycles. The audit findings are recorded as a milestone-level synchronization checkpoint, and remediation issues are scoped and prioritized before the next recursive phase begins.
